\chapter{Боковой теплообмен по закону Ньютона}

\section{Постановка практического задания 4}

В пункте 4 практической части требуется решить вторую практическую
задачу с учётом теплообмена на боковой поверхности. Поэтому сохраняются
начальная температура $u(x,0)=x$ и постоянные индивидуальные
неоднородности $A,B$ на торцах. Гармонические данные из практического
задания 3 здесь не используются.

Для варианта 11 типы условий: первый род при $x=0$ и второй род
при $x=l$. Сохраняя знак общей записи лабораторной работы,
получаем
\[
 u(0,t)=A,\qquad -u_x(l,t)=B.
\]
Здесь $B$ обозначает отрицательную координатную производную
температуры. Если вместо этого определить $B=u_x(l,t)$, во всех
последующих формулах следует заменить $B$ на $-B$.

Обозначим постоянную температуру среды через $T_e$.
Для однородного стержня с площадью сечения $S$ и периметром сечения $P$
закон Ньютона на боковой поверхности даёт уравнение
\[
 cS u_t=kS u_{xx}-HP(u-T_e).
\]
Здесь $c$ --- объёмная теплоёмкость, $k$ --- теплопроводность,
$H>0$ --- коэффициент теплоотдачи. Внутренние источники отсутствуют.
Введём
\[
 a^2=\frac{k}{c},\qquad b=\frac{HP}{cS}>0.
\]
Для круглого сечения радиуса $R$ имеем $b=2H/(cR)$.
Таким образом,
\[
\boxed{\begin{cases}
 u_t=a^2u_{xx}-b(u-T_e),&0<x<l,\quad t>0,\\
 u(0,t)=A,\quad -u_x(l,t)=B,&t>0,\\
 u(x,0)=x,&0\le x\le l.
\end{cases}}
\]
В обозначениях $u_t=a^2u_{xx}-bu+f$ функция $f=bT_e$ возникает
из теплообмена, а не из дополнительного внутреннего источника.

\section{Стационарный профиль при постоянной среде}

Пусть $u(x,t)=U_{T_e}(x)$. Тогда
\[
 a^2U_{T_e}''-b(U_{T_e}-T_e)=0,\qquad
 U_{T_e}(0)=A,\quad U_{T_e}'(l)=-B.
\]
Обозначим
\[
 \kappa=\frac{\sqrt b}{a}>0.
\]
Положим $Y=U_{T_e}-T_e$. Тогда
\[
 Y''-\kappa^2Y=0,\qquad Y(0)=A-T_e,\quad Y'(l)=-B.
\]
Удобно представить общее решение относительно правого торца:
\[
 Y(x)=C\cosh(\kappa(l-x))+D\sinh(\kappa(l-x)).
\]
Из $Y'(l)=-B$ получаем $D=B/\kappa$, а из левого условия
\[
 C=\frac{A-T_e-(B/\kappa)\sinh(\kappa l)}{\cosh(\kappa l)}.
\]
Применяя формулу разности гиперболических синусов, получаем
\[
\boxed{U_{T_e}(x)=T_e+
 (A-T_e)\frac{\cosh(\kappa(l-x))}{\cosh(\kappa l)}
 -\frac{B}{\kappa}\frac{\sinh(\kappa x)}{\cosh(\kappa l)}.}
\]
Непосредственная проверка даёт
\[
 U_{T_e}(0)=A,\qquad U_{T_e}'(l)=-B,\qquad
 U_{T_e}''=\kappa^2(U_{T_e}-T_e).
\]
Стационарная задача имеет единственное решение. Действительно,
для разности $W$ двух решений имеем $a^2W''-bW=0$,
$W(0)=0$, $W'(l)=0$. Умножая уравнение на $W$ и интегрируя,
получаем
\[
 -a^2\int_0^l(W')^2\,dx-b\int_0^lW^2\,dx=0,
\]
откуда $W\equiv0$.

\section{Полное решение методом Фурье}

Представим температуру как сумму стационарной и переходной частей:
\[
 u(x,t)=U_{T_e}(x)+v(x,t).
\]
Тогда
\[
\begin{cases}
 v_t=a^2v_{xx}-bv,\\
 v(0,t)=0,\quad v_x(l,t)=0,\\
 v(x,0)=x-U_{T_e}(x).
\end{cases}
\]
Собственные функции и скорости затухания:
\[
 X_n(x)=\sin(k_nx),\quad
 k_n=\frac{(n+\tfrac12)\pi}{l},\quad
 \sigma_n=a^2k_n^2+b,\qquad n=0,1,2,\ldots.
\]
Следовательно,
\[
\boxed{u(x,t)=U_{T_e}(x)+\sum_{n=0}^{\infty}
 C_n(T_e)e^{-\sigma_nt}\sin(k_nx),}
\]
где
\[
\boxed{C_n(T_e)=\frac2l\int_0^l
 \bigl(\xi-U_{T_e}(\xi)\bigr)\sin(k_n\xi)\,d\xi.}
\]
Это уже явная формула решения. Можно вычислить коэффициенты
без интегралов. Поскольку $\cos(k_nl)=0$, $\sin(k_nl)=(-1)^n$,
\[
 \int_0^l\xi\sin(k_n\xi)\,d\xi=\frac{(-1)^n}{k_n^2}.
\]
Умножим равенство $U_{T_e}''=\kappa^2(U_{T_e}-T_e)$
на $\sin(k_nx)$ и дважды проинтегрируем по частям.
Граничные члены дают
\[
 (k_n^2+\kappa^2)\int_0^lU_{T_e}(x)\sin(k_nx)\,dx
 =Ak_n-B(-1)^n+\frac{\kappa^2T_e}{k_n}.
\]
Таким образом,
\[
\boxed{C_n(T_e)=\frac2l\left[
 \frac{(-1)^n}{k_n^2}
 -\frac{Ak_n-B(-1)^n+\kappa^2T_e/k_n}{k_n^2+\kappa^2}
 \right].}
\]
Для произвольной начальной функции $\phi$ вместо $x$
в интегральной формуле используется $\phi(\xi)$.
При $\phi=x$ согласование с границами в момент $t=0$
возможно лишь при $A=0$, $B=-1$.
Для остальных значений начальное условие понимается в смысле
$L^2(0,l)$, а решение гладко при $t>0$.
Даже при согласовании $\phi=x$ не является стационарной
температурой при $b>0$ и постоянном $T_e$.

\section{Два случая температуры среды}

\subsection{Температура среды равна нулю}

При $T_e=0$ уравнение имеет вид
\[
 u_t=a^2u_{xx}-bu.
\]
Стационарная температура равна
\[
\boxed{U_0(x)=A\frac{\cosh(\kappa(l-x))}{\cosh(\kappa l)}
 -\frac B\kappa\frac{\sinh(\kappa x)}{\cosh(\kappa l)}.}
\]
Полное решение:
\[
 u(x,t)=U_0(x)+\sum_{n=0}^{\infty}
 C_n(0)e^{-(a^2k_n^2+b)t}\sin(k_nx),
\]
\[
 C_n(0)=\frac2l\left[
 \frac{(-1)^n}{k_n^2}
 -\frac{Ak_n-B(-1)^n}{k_n^2+\kappa^2}\right].
\]
Если $A=B=0$, то $U_0\equiv0$ и
\[
 u(x,t)=\frac2l\sum_{n=0}^{\infty}\frac{(-1)^n}{k_n^2}
 e^{-(a^2k_n^2+b)t}\sin(k_nx)\longrightarrow0.
\]
Если хотя бы одна граничная неоднородность ненулевая,
стационарный профиль ненулевой: торцы поддерживают температуру
или поток, компенсирующие боковые потери.

\subsection{Температура среды постоянна}

При $T_e=T_0=\mathrm{const}$ получаем
\[
 u_t=a^2u_{xx}-bu+bT_0.
\]
Стационарная температура и полное решение:
\[
\boxed{U_{T_0}(x)=T_0+
 (A-T_0)\frac{\cosh(\kappa(l-x))}{\cosh(\kappa l)}
 -\frac B\kappa\frac{\sinh(\kappa x)}{\cosh(\kappa l)},}
\]
\[
 u(x,t)=U_{T_0}(x)+\sum_{n=0}^{\infty}
 C_n(T_0)e^{-(a^2k_n^2+b)t}\sin(k_nx).
\]
При одинаковых $A,B$ отличие от случая нулевой среды равно
\[
\boxed{U_{T_0}(x)-U_0(x)=T_0\left[
 1-\frac{\cosh(\kappa(l-x))}{\cosh(\kappa l)}\right].}
\]
На левом торце разность нулевая: его температура фиксирована
и равна $A$. Поэтому нельзя просто прибавить $T_0$ к $U_0$,
не изменив одновременно левое граничное значение.

Если $A=T_0$ и $B=0$, то
\[
 U_{T_0}(x)\equiv T_0.
\]
Это необходимое и достаточное условие пространственно постоянного
стационарного профиля: постоянство требует $U=T_0$
из уравнения и $A=T_0$, $B=0$ из граничных условий.
Если $A=B=0$, но $T_0\ne0$, профиль ненулевой:
\[
 U_{T_0}(x)=T_0\left[
 1-\frac{\cosh(\kappa(l-x))}{\cosh(\kappa l)}\right].
\]
При $b>0$ тождественно нулевой стационарный режим возможен
тогда и только тогда, когда $A=B=T_0=0$.

\section{Асимптотика и проверка предела}

Все скорости затухания положительны:
\[
 \sigma_n=a^2k_n^2+b>0.
\]
Поэтому в обоих случаях
\[
\boxed{u(\cdot,t)\longrightarrow U_{T_e}\quad(t\to\infty).}
\]
Начальное распределение влияет только на переходный процесс.
В норме $L^2(0,l)$ справедлива оценка
\[
 \|u(\cdot,t)-U_{T_e}\|_{L^2}
 \le e^{-(b+a^2\pi^2/(4l^2))t}
 \|x-U_{T_e}\|_{L^2}.
\]
Сходимость равномерна по $x$ при $t\to\infty$ вследствие
сглаживания решения и экспоненциального затухания ряда.
Боковой теплообмен увеличивает скорости затухания на $b$.
При этом стационарный профиль зависит от $b$ через $\kappa$.

Полезная проверка: при исчезновении теплообмена, $b\downarrow0$,
\[
 \frac{\cosh(\kappa(l-x))}{\cosh(\kappa l)}\longrightarrow1,
 \qquad
 \frac{\sinh(\kappa x)}{\kappa\cosh(\kappa l)}\longrightarrow x.
\]
Поэтому
\[
\boxed{U_{T_e}(x)\longrightarrow A-Bx.}
\]
Температура среды исчезает из предельной формулы, как и должно быть
при отсутствии бокового теплообмена. Это совпадает со стационарным
режимом второй практической задачи.
